\documentclass[a4paper,11pt,norsk]{article}
\usepackage{packages}

\begin{document}

\input{Header/overskrift}

\section*{Oppgave 1.}
\begin{enumerate}
    \item 
        \[
            \Delta \lambda = \frac{2v f_0}{c} \approx \SI{1.67}{\kilo\hertz}
        \]
    \item Ren dopplerradar sender kunn ut et kontinuerlig signal med en fast frekvens, der den
        bruker faseforskjellen mellom utsendt og mottatt signal for å måle fart. Det er ignenting med signalet 
        som sier oss noe om hvor langt signalet beveget seg før den ble reflektert, med mindre man vet godt hvordan tingen
        man måler ser ut ifht rf reflektivitet, da man muligens kan se på styrken til det mottatte signalet (ren spekulasjon fra min side).
    \item Blir kosinusdelen av geometrien til oppsettet, så $\cos(45^\circ) \approx 0.707$ av det vi fikk i (a),
        så omtrent $\SI{1.18}{\kilo\hertz}$.
\end{enumerate}

\section*{Oppgave 2.}
\begin{enumerate}
    \item Radarligningen
        \[
            P_r = \frac{P_t G^2 \lambda^2 \sigma}{(4\pi)^3 R^4}
        \]
        Vi har $G = 28 \text{ dBi } = 631$, $\lambda = \frac{c}{f} = \frac{3 \cdot 10^8}{3.1 \cdot 10^9} \approx \SI{0.0968}{\meter}$. 
        Løser for $R$ for å få maks radius 
        \[
            R^4 = \frac{10^4 \cdot 631^2 \cdot 0.0968^2 \cdot 10}{(4\pi)^3 \cdot 4 \cdot 10^{-10}} \approx
        \]
        det gir
        \[
            R_{max} \approx \SI{4656}{\meter}
        \]
    \item 
        Da får vi
        \[
            R_{max} \approx \SI{5537}{\meter}
        \]
\end{enumerate}

\section*{Oppgave 3.}
\begin{enumerate}
    \item Vi har
        \[
            A_{eff} = \frac{a^2}{\sqrt{3}}
        \]
        og får da radartverrsnitt
        \[
            \sigma = \frac{4\pi A_{eff}^2}{\lambda^2} = \frac{4\pi a^4}{3\lambda^2} \approx \SI{16.4}{\meter^2}
        \]
    \item $\sigma = \pi r^2 \implies r = \sqrt{\frac{16.4}{\pi}} = \SI{2.28}{\meter}$
\end{enumerate}

\end{document}
